%% ============================================================
%% Introduction to Cognitive Psychology — Week 6 Study Notes
%% ============================================================
\documentclass[11pt,a4paper]{article}

\usepackage[a4paper, left=1.8cm, right=1.8cm, top=1.3cm, bottom=1.8cm]{geometry}
\usepackage{booktabs}
\usepackage{tabularx}
\usepackage{array}
\usepackage{enumitem}
\usepackage{titlesec}
\usepackage{fancyhdr}
\usepackage{fontenc}
\usepackage{inputenc}
\usepackage{microtype}
\usepackage{parskip}
\usepackage{hyperref}
\usepackage{amsmath}
\usepackage{amssymb}

\hypersetup{hidelinks, colorlinks=false}

%% ── Table helpers ─────────────────────────────────────────
\newcolumntype{L}[1]{>{\raggedright\arraybackslash}p{#1}}

%% ── Section headings ──────────────────────────────────────
\titleformat{\section}{\large\bfseries}{}{0em}{}[\titlerule]
\titleformat{\subsection}{\normalsize\bfseries}{}{0em}{}
\titleformat{\subsubsection}{\normalsize\bfseries}{}{0em}{}

\titlespacing{\section}{0pt}{12pt}{4pt}
\titlespacing{\subsection}{0pt}{8pt}{3pt}
\titlespacing{\subsubsection}{0pt}{6pt}{2pt}

%% ── Page style ────────────────────────────────────────────
\pagestyle{fancy}
\fancyhf{}
\renewcommand{\headrulewidth}{0pt}
\fancyfoot[C]{\small Cognitive Psychology --- Week 6 \textbar\ Page \thepage}

%% ── Misc helpers ──────────────────────────────────────────
\newcommand{\bb}[1]{\textbf{#1}}
\setlength{\parskip}{4pt}

%% ═══════════════════════════════════════════════════════════
\begin{document}
\setlength{\arrayrulewidth}{0.4pt}

\begin{center}
\bfseries\LARGE Introduction to Cognitive Psychology\par
\vspace{0.2cm}
\large Assignment 6 Detailed Solutions
\end{center}
\vspace{0.5cm}

%% ════════════════════════════════════════════════════════════
\section*{ASSIGNMENT 6 --- Full Solutions with Explanations}

\subsubsection*{Question 1}
\textit{According to the notion of cognitive economy, a characteristic like "has wings" would be stored along with which of the following semantic memory nodes.}
\medskip

\begin{tabular}{L{0.3cm} L{14.9cm}}
& \textbf{A.\ Bird \checkmark} \\
& B.\ Ostrich \\
& C.\ Hummingbird \\
& D.\ All of the above \\
\end{tabular}

\vspace{0.3cm}
\noindent\textbf{Ans:} A — Bird\\[0.2cm]
\textbf{Why this answer:} \\
According to \textbf{cognitive economy}, properties shared by an entire category are stored at the \textbf{highest applicable level} in the conceptual hierarchy rather than being redundantly repeated at every subordinate node. Since "has wings" is a property common to virtually all birds, it is stored at the \textbf{Bird} node — not repeated at subordinate nodes like Ostrich or Hummingbird. This principle was central to Collins \& Quillian's (1969) \textbf{hierarchical network model} of semantic memory.
\vspace{0.4cm}

\subsubsection*{Question 2}
\textit{Priming in lexical decision tasks may be explained by the idea of:}
\medskip

\begin{tabular}{L{0.3cm} L{14.9cm}}
& A.\ Episodic memory \\
& B.\ Encoding specificity \\
& \textbf{C.\ Spreading activation \checkmark} \\
& D.\ Typicality effects \\
\end{tabular}

\vspace{0.3cm}
\noindent\textbf{Ans:} C — Spreading activation\\[0.2cm]
\textbf{Why this answer:} \\
In a \textbf{lexical decision task}, participants decide whether a letter string is a real word. \textbf{Priming} occurs when a preceding related word (the prime) speeds up responses to a target word. This is best explained by \textbf{spreading activation} (Collins \& Loftus, 1975): activation from the prime node spreads through the semantic network to related nodes, pre-activating them and lowering the threshold needed for recognition. For example, "bread" primes "butter" because the two nodes are closely linked in the network.
\vspace{0.4cm}

\subsubsection*{Question 3}
\textit{Which of the following is a good example of a superordinate level of categorization?}
\medskip

\begin{tabular}{L{0.3cm} L{14.9cm}}
& \textbf{A.\ Fruit \checkmark} \\
& B.\ Banana \\
& C.\ Fuji apple \\
& D.\ Golden delicious apple \\
\end{tabular}

\vspace{0.3cm}
\noindent\textbf{Ans:} A — Fruit\\[0.2cm]
\textbf{Why this answer:} \\
Rosch's (1978) \textbf{basic-level theory} identifies three levels in the conceptual hierarchy: \textbf{superordinate} (most general), \textbf{basic}, and \textbf{subordinate} (most specific). In the domain of food: *Fruit* is the \textbf{superordinate} level, *Banana* or *Apple* is the \textbf{basic} level, and *Fuji apple* or *Golden Delicious apple* is the \textbf{subordinate} level. Superordinate categories are the broadest and share the fewest features across members.
\vspace{0.4cm}

\subsubsection*{Question 4}
\textit{The schema view of concept formation assumes that:}
\medskip

\begin{tabular}{L{0.3cm} L{14.9cm}}
& A.\ There are clear boundaries among individual schemata \\
& B.\ There is cognitive economy among concepts \\
& \textbf{C.\ Information is abstracted across instances \checkmark} \\
& D.\ No information is stored about actual instances \\
\end{tabular}

\vspace{0.3cm}
\noindent\textbf{Ans:} C — Information is abstracted across instances\\[0.2cm]
\textbf{Why this answer:} \\
The \textbf{schema view} of concept formation holds that people extract and store \textbf{abstract, generalized representations} formed across many individual encounters with a category. Rather than storing a perfect prototype or every individual exemplar, schemata capture the common structure, typical attributes, and typical relationships found across instances. Note that schemata do \textbf{not} require perfectly clear boundaries (option A) and do \textbf{not} claim that no instance-specific information is ever retained (option D).
\vspace{0.4cm}

\subsubsection*{Question 5}
\textit{According to research by Collins and Quillian, the statement "Siamese cats have blue eyes" will be verified:}
\medskip

\begin{tabular}{L{0.3cm} L{14.9cm}}
& A.\ Slower than "Siamese cats give birth to live young" \\
& \textbf{B.\ Faster than "Siamese cats give birth to live young" \checkmark} \\
& C.\ In the same amount of time as "Siamese cats give birth to live young" \\
& D.\ Slower than "Siamese cats have tails" \\
\end{tabular}

\vspace{0.3cm}
\noindent\textbf{Ans:} B — Faster than "Siamese cats give birth to live young"\\[0.2cm]
\textbf{Why this answer:} \\
Collins \& Quillian's \textbf{hierarchical network model} predicts that verification time increases with the number of hierarchical levels that must be traversed. "Siamese cats have blue eyes" is a \textbf{subordinate-level} property stored directly at the *Siamese cat* node, requiring minimal traversal. By contrast, "Siamese cats give birth to live young" is a \textbf{mammal-level} property, requiring the system to travel up through *Cat → Mammal*, taking longer. Therefore, the blue-eyes statement is verified \textbf{faster}.
\vspace{0.4cm}

\subsubsection*{Question 6}
\textit{Conrad has found evidence that the statement "A shark can move" can be verified in the same amount of time as "An animal can move." These results suggest that reaction time is best predicted by:}
\medskip

\begin{tabular}{L{0.3cm} L{14.9cm}}
& A.\ Cognitive economy \\
& \textbf{B.\ Frequency of associations \checkmark} \\
& C.\ Encoding specificity \\
& D.\ Episodic memory \\
\end{tabular}

\vspace{0.3cm}
\noindent\textbf{Ans:} B — Frequency of associations\\[0.2cm]
\textbf{Why this answer:} \\
Collins \& Quillian's hierarchical model predicts that "A shark can move" (two levels up) should be verified \textbf{slower} than "An animal can move" (zero levels for *Animal*). Conrad's finding that both statements are verified equally fast challenges the strict hierarchical distance account. Instead, it supports the idea that \textbf{associative strength (frequency of association)} — how often "shark" and "move" co-occur in experience — is a better predictor of reaction time than hierarchical distance alone. This was a key insight that motivated Collins \& Loftus's later \textbf{spreading activation} model.
\vspace{0.4cm}
\newpage

\subsubsection*{Question 7}
\textit{Which of the following would be a part of your declarative memory system?}
\medskip

\begin{tabular}{L{0.3cm} L{14.9cm}}
& A.\ Knowing how to ride a bicycle \\
& B.\ Knowing how to drive a car \\
& C.\ Knowing how to react at a red light \\
& \textbf{D.\ Being able to name a hybrid car \checkmark} \\
\end{tabular}

\vspace{0.3cm}
\noindent\textbf{Ans:} D — Being able to name a hybrid car\\[0.2cm]
\textbf{Why this answer:} \\
\textbf{Declarative memory} (also called explicit memory) encompasses knowledge that can be \textbf{consciously recalled and verbally stated}, including both \textbf{episodic} (personal events) and \textbf{semantic} (general factual knowledge) memory. Being able to *name* a hybrid car is accessing a \textbf{semantic fact} — a classic example of declarative knowledge. The other options (riding a bicycle, driving a car, reacting at a red light) all involve \textbf{procedural memory}, a form of \textbf{non-declarative (implicit) memory} that encodes motor skills and habits not easily verbalized.
\vspace{0.4cm}

\subsubsection*{Question 8}
\textit{Collins and Loftus's spreading activation theory differs from the hierarchical network theory in that:}
\medskip

\begin{tabular}{L{0.3cm} L{14.9cm}}
& \textbf{A.\ It dispenses with the idea of cognitive economy \checkmark} \\
& B.\ It relies on the assumption of hierarchical structure \\
& C.\ It cannot account for the typicality effect \\
& D.\ It makes stronger predictions than hierarchical models \\
\end{tabular}

\vspace{0.3cm}
\noindent\textbf{Ans:} A — It dispenses with the idea of cognitive economy\\[0.2cm]
\textbf{Why this answer:} \\
Collins \& Quillian's original hierarchical model relied on \textbf{cognitive economy} — properties are stored at the highest applicable level to avoid redundancy. Collins \& Loftus (1975) revised this in their \textbf{spreading activation model}: concepts are organized by \textbf{semantic relatedness} (not rigid hierarchy), and \textbf{properties can be stored at multiple nodes} regardless of hierarchical level. This eliminates strict cognitive economy. The spreading activation model also better accounts for \textbf{typicality effects} (option C), \textbf{associative strength}, and other phenomena that the hierarchical model failed to explain.
\vspace{0.4cm}

\subsubsection*{Question 9}
\textit{"Characteristic features" and "family resemblance" are important aspects of the \_\_\_\_\_\_\_\_ view of concepts.}
\medskip

\begin{tabular}{L{0.3cm} L{14.9cm}}
& A.\ Classical \\
& \textbf{B.\ Prototype \checkmark} \\
& C.\ Exemplar \\
& D.\ Schema \\
\end{tabular}

\vspace{0.3cm}
\noindent\textbf{Ans:} B — Prototype\\[0.2cm]
\textbf{Why this answer:} \\
The \textbf{prototype view} (Rosch, 1973) holds that people represent categories as a \textbf{prototype} — an idealized, average representation capturing the most \textbf{characteristic features} of category members. Members are judged as category instances based on their resemblance to this prototype. \textbf{Family resemblance} (Wittgenstein) captures the idea that members share overlapping features without any single feature being necessary or sufficient — exactly the logic underlying the prototype view. The \textbf{classical view} requires necessary and sufficient features (clear boundaries), while \textbf{exemplar} views store actual encountered instances, and \textbf{schema} views emphasize abstracted relational structures.
\vspace{0.4cm}
\newpage
\subsubsection*{Question 10}
\textit{Studies of semantic memory have shown that in a lexical decision task, people are faster at responding to the stimulus "bread" if it is paired with a stimulus such as:}
\medskip

\begin{tabular}{L{0.3cm} L{14.9cm}}
& A.\ "rencle" \\
& B.\ "dog" \\
& C.\ "island" \\
& \textbf{D.\ "butter" \checkmark} \\
\end{tabular}

\vspace{0.3cm}
\noindent\textbf{Ans:} D — "butter"\\[0.2cm]
\textbf{Why this answer:} \\
This is a classic demonstration of \textbf{semantic priming} via \textbf{spreading activation}. When the word "bread" is processed, activation spreads through the semantic network to closely associated nodes — and "butter" is one of the strongest associates of "bread" (they frequently co-occur and are semantically/functionally related). This pre-activation lowers the threshold for recognizing "butter" as a real word, resulting in a \textbf{faster lexical decision}. The other options ("rencle" is a non-word, "dog" and "island" are semantically unrelated) would produce no such priming benefit.
\vspace{0.4cm}

\medskip
\subsection*{Assignment Quick Answer Summary}

\begin{center}
\renewcommand{\arraystretch}{1.4}
\begin{tabular}{|L{0.8cm}|L{6.2cm}|L{8.2cm}|}
\hline
\bb{Q\#} & \bb{Topic} & \bb{Correct Answer} \\
\hline
1 & Cognitive Economy & A — Bird \\
2 & Priming in Lexical Decision & C — Spreading activation \\
3 & Superordinate Categorization & A — Fruit \\
4 & Schema View of Concepts & C — Information abstracted \\
5 & Collins \& Quillian Verification & B — Faster \\
6 & Conrad's Findings & B — Frequency of associations \\
7 & Declarative Memory & D — Naming a hybrid car \\
8 & Spreading Activation vs. Hierarchy & A — Dispenses with cognitive economy \\
9 & Prototype View & B — Prototype \\
10 & Semantic Priming & D — "butter" \\
\hline
\end{tabular}
\end{center}

\medskip
\end{document}
